\chapter{Chapter06 — Low Degree Junta Learning Algorithm}
%%% generated by the blueprint pipeline from FABL.Chapter06.LearningAndTesting.LowDegreeJuntaLearningAlgorithm
\section{Overview}

Declarations of the Lean module \texttt{FABL.Chapter06.LearningAndTesting.LowDegreeJuntaLearningAlgorithm}.

\section{Declarations}

\begin{definition}[Low degree junta learning failure parameter]
\lean{FABL.lowDegreeJuntaLearningFailure}
\statementsource{Analysis-of-Boolean-Functions-by-Ryan-ODonnell}{
  pdf-ffbe9c7b3306-p079-b001,
  pdf-ffbe9c7b3306-p081-b003,
  pdf-ffbe9c7b3306-p156-b006,
  pdf-ffbe9c7b3306-p156-b007,
  pdf-ffbe9c7b3306-p156-b008
}
\label{FABL.lowDegreeJuntaLearningFailure}
\leanok
\uses{FABL.PositiveLearningParameter}
The low degree junta learning failure parameter is the rational number $1/16$, regarded
as a learning parameter in the interval $(0,1/2]$.
\end{definition}

\begin{definition}[Decision tree from an optional output]
\lean{FABL.lowDegreeJuntaHypothesisOfOutput}
\label{FABL.lowDegreeJuntaHypothesisOfOutput}
\leanok
\uses{FABL.DecisionTree, FABL.DecisionTreeHypothesis, FABL.DecisionTreeHypothesis.ofF₂DecisionTree, FABL.Sign}
For a natural number $n$, associate to an optional decision tree on the $n$ coordinates
a decision tree hypothesis with leaves labeled by signs in $\{\pm1\}$. If no tree is
supplied, the hypothesis is the constant leaf labeled $+1$; if a tree is supplied,
retain its query coordinates and recursively convert its two child trees.
\end{definition}

\begin{definition}[Low degree junta learning program]
\lean{FABL.lowDegreeJuntaLearningProgram}
\statementsource{Analysis-of-Boolean-Functions-by-Ryan-ODonnell}{
  pdf-ffbe9c7b3306-p081-b003,
  pdf-ffbe9c7b3306-p081-b004
}
\label{FABL.lowDegreeJuntaLearningProgram}
\leanok
\uses{FABL.DecisionTreeHypothesis, FABL.LearningAccuracy, FABL.LearningProgram, FABL.LearningProgram.map, FABL.lowDegreeJuntaHypothesisOfOutput, FABL.lowDegreeJuntaLearnerProgram, FABL.lowDegreeJuntaLearningFailure}
For natural numbers $n$ and $k$ and a rational accuracy parameter in $[0,1/2]$, define a
random-example learning program whose output is a decision tree on the $n$ coordinates.
The program applies the specified conversion from the low-degree junta learner’s output
to a decision-tree hypothesis, using parameters $n$ and $k$; the accuracy parameter does
not affect the program.
\end{definition}

\begin{definition}[Class of k-junta Boolean functions]
\lean{FABL.kJuntaConceptClass}
\statementsource{Analysis-of-Boolean-Functions-by-Ryan-ODonnell}{
  pdf-ffbe9c7b3306-p044-b003,
  pdf-ffbe9c7b3306-p155-b010,
  pdf-ffbe9c7b3306-p168-b001,
  pdf-ffbe9c7b3306-p173-b002
}
\label{FABL.kJuntaConceptClass}
\leanok
\uses{FABL.BooleanFunction, FABL.IsKJunta}
For natural numbers $n$ and $k$, the $k$-junta concept class on the $n$-dimensional sign
cube is the set of Boolean functions on that cube that are $k$-juntas.
\end{definition}

\begin{theorem}[Evaluation of the converted decision tree]
\lean{FABL.evaluate_lowDegreeJuntaHypothesisOfOutput_eq}
\label{FABL.evaluate_lowDegreeJuntaHypothesisOfOutput_eq}
\leanok
\uses{FABL.BooleanFunction, FABL.DecisionTree, FABL.DecisionTreeComputesTarget, FABL.DecisionTreeHypothesis.eval, FABL.DecisionTreeHypothesis.eval_ofF₂DecisionTree, FABL.DecisionTreeHypothesis.evaluate, FABL.DecisionTreeHypothesis.ofF₂DecisionTree, FABL.Sign, FABL.binaryCubeSignEquiv, FABL.lowDegreeJuntaHypothesisOfOutput}
\difficulty{3}
Let $n$ be a nonnegative integer, let $f$ be a function from the $n$-dimensional binary
cube to the signs, and let $T$ be a decision tree on that cube with sign outputs. If $T$
computes $f$ after identifying binary coordinates with sign coordinates, then the
sign-valued hypothesis tree obtained from $T$ evaluates to $f$ on every sign vector.
\end{theorem}

\begin{theorem}[Accuracy of the output hypothesis]
\lean{FABL.lowDegreeJuntaHypothesisOfOutput_accurate_of_not_bad}
\label{FABL.lowDegreeJuntaHypothesisOfOutput_accurate_of_not_bad}
\leanok
\uses{FABL.BooleanFunction, FABL.DecisionTree, FABL.DecisionTreeComputesTarget, FABL.DecisionTreeHypothesis.evaluate, FABL.JuntaLearnerOutputBad, FABL.LearningAccuracy, FABL.Sign, FABL.evaluate_lowDegreeJuntaHypothesisOfOutput_eq, FABL.lowDegreeJuntaHypothesisOfOutput, FABL.relativeHammingDist}
\difficulty{4}
Let $n$ be a natural number, let $f:\{−1,1\}^n\to\{−1,1\}$ be a Boolean function, let
$a\in[0,1/2]\cap\mathbb{Q}$, and let the learner's output be either absent or a decision
tree on the $n$ binary coordinates with values in $\{−1,1\}$. If the output is not
bad—that is, it is present and its tree computes $f$ on every input—then the hypothesis
obtained from the output has relative Hamming distance at most $a$ from $f$.
\end{theorem}

\begin{theorem}[Probability bound from an output cover]
\lean{FABL.one_le_eventProbability_add_of_output_cover}
\label{FABL.one_le_eventProbability_add_of_output_cover}
\leanok
\uses{FABL.BooleanFunction, FABL.LearningAccess, FABL.LearningCost, FABL.LearningProgram, FABL.LearningProgram.eventProbability, FABL.LearningProgram.runWithCost, FABL.pmfToOuterMeasure_ne_top}
\difficulty{5}
Let a learning program with any access mode have outputs in a type $\alpha$, and let its
target be a Boolean function. If every possible output satisfies at least one of two
properties $A$ and $B$, then the probability that the program's output satisfies $A$
plus the probability that its output satisfies $B$ is at least $1$.
\end{theorem}

\begin{definition}[Low degree junta learning algorithm]
\lean{FABL.lowDegreeJuntaLearningAlgorithm}
\statementsource{Analysis-of-Boolean-Functions-by-Ryan-ODonnell}{
  pdf-ffbe9c7b3306-p081-b003,
  pdf-ffbe9c7b3306-p156-b002,
  pdf-ffbe9c7b3306-p156-b004,
  pdf-ffbe9c7b3306-p156-b005,
  pdf-ffbe9c7b3306-p156-b006,
  pdf-ffbe9c7b3306-p156-b007,
  pdf-ffbe9c7b3306-p156-b008
}
\label{FABL.lowDegreeJuntaLearningAlgorithm}
\leanok
\uses{FABL.DecisionTreeHypothesis.evaluate, FABL.DecisionTreeHypothesis.finiteRepresentation, FABL.LearningAlgorithm, FABL.LearningProgram.eventProbability, FABL.LearningProgram.runWithCost_map, FABL.juntaTreeCallCount, FABL.juntaTreePerCallFailure, FABL.lowDegreeJuntaLearnerProgram_cost_le, FABL.lowDegreeJuntaLearningFailure, FABL.lowDegreeJuntaLearningProgram, FABL.lowDegreeJuntaNodeRandomExampleBound, FABL.lowDegreeJuntaNodeUniformWorkBound, FABL.relativeHammingDist}
For natural numbers $n$ and $k$, define a random-example learning algorithm whose
program at each requested accuracy is the low degree junta learning program with
parameters $n$ and $k$. Its success probability is the probability that the output
decision tree has relative Hamming distance at most the requested accuracy from the
target; its uniform bounds on random-example calls and local work are, respectively, the
junta tree call count times the node random-example bound and the junta tree call count
times the node uniform-work bound, both evaluated with the per-call failure parameter
derived from $k$ and $1/16$. Its membership-query bound is zero.
\end{definition}

\begin{theorem}[Low degree junta learning algorithm success probability ge]
\lean{FABL.lowDegreeJuntaLearningAlgorithm_successProbability_ge}
\statementsource{Analysis-of-Boolean-Functions-by-Ryan-ODonnell}{
  pdf-ffbe9c7b3306-p156-b005,
  pdf-ffbe9c7b3306-p156-b006,
  pdf-ffbe9c7b3306-p156-b008,
  pdf-ffbe9c7b3306-p081-b003,
  pdf-ffbe9c7b3306-p081-b004
}
\label{FABL.lowDegreeJuntaLearningAlgorithm_successProbability_ge}
\leanok
\uses{FABL.BooleanFunction, FABL.DecisionTree, FABL.DecisionTreeHypothesis.evaluate, FABL.IsKJunta, FABL.JuntaLearnerOutputBad, FABL.LearningAccuracy, FABL.LearningProgram.eventProbability, FABL.LearningProgram.map, FABL.Sign, FABL.eventProbability_map_output_eq, FABL.lowDegreeJuntaHypothesisOfOutput, FABL.lowDegreeJuntaHypothesisOfOutput_accurate_of_not_bad, FABL.lowDegreeJuntaLearnerProgram, FABL.lowDegreeJuntaLearnerProgram_failureProbability_le, FABL.lowDegreeJuntaLearningAlgorithm, FABL.lowDegreeJuntaLearningFailure, FABL.lowDegreeJuntaLearningProgram, FABL.one_le_eventProbability_add_of_output_cover, FABL.relativeHammingDist}
\difficulty{5}
% dataset-pipeline: stub -- prose pending build_informal_statements.py
Statement pending regeneration.
\end{theorem}

\begin{theorem}[Low degree junta learning algorithm learns]
\lean{FABL.lowDegreeJuntaLearningAlgorithm_learns}
\statementsource{Analysis-of-Boolean-Functions-by-Ryan-ODonnell}{
  pdf-ffbe9c7b3306-p156-b002,
  pdf-ffbe9c7b3306-p081-b003,
  pdf-ffbe9c7b3306-p081-b004
}
\label{FABL.lowDegreeJuntaLearningAlgorithm_learns}
\leanok
\uses{FABL.LearningAccuracy, FABL.LearnsConceptClassWithError, FABL.kJuntaConceptClass, FABL.lowDegreeJuntaLearningAlgorithm, FABL.lowDegreeJuntaLearningAlgorithm_successProbability_ge}
\difficulty{1}
% dataset-pipeline: stub -- prose pending build_informal_statements.py
Statement pending regeneration.
\end{theorem}
